\documentclass[10pt,aspectratio=169]{beamer}

%-------------------------------------------------
% Theme and packages
%-------------------------------------------------
\usetheme{Madrid}
\usecolortheme{seahorse}
\useoutertheme{infolines}
\useinnertheme{rounded}

\usepackage[utf8]{inputenc}
\usepackage[T1]{fontenc}
\usepackage{amsmath,amssymb}
\usepackage{graphicx}
\usepackage{xcolor}
\usepackage{hyperref}

% Slightly tighter list spacing
\setbeamertemplate{itemize items}[circle]
\setbeamertemplate{itemize subitem}[triangle]
\setbeamertemplate{navigation symbols}{}

% Highlight box for reviewer quotes
\setbeamercolor{quotebox}{bg=blue!8,fg=black}
\newcommand{\reviewerquote}[1]{%
  \begin{beamercolorbox}[sep=6pt,rounded=true,shadow=true]{quotebox}%
  \itshape #1%
  \end{beamercolorbox}%
}

%-------------------------------------------------
% Title information
%-------------------------------------------------
\title[SFF Evaluation -- Interpretation]{Interpretations of the SFF Evaluation \\ \small Center for Quantum Control (Project 362317)}
\subtitle{From Phase 1 (Mark 6) toward an Outstanding Phase 2 Proposal}
\author{Morten Hjorth-Jensen}
\institute[UiO]{University of Oslo}
\date{\today}

\begin{document}

%-------------------------------------------------
\begin{frame}
  \titlepage
\end{frame}

%-------------------------------------------------
\begin{frame}{Outline}
  \tableofcontents
\end{frame}

%=================================================
\section{Overall impression}
%=================================================
\begin{frame}{Overall impression of the Phase 1 evaluation}
  \begin{itemize}
    \item The review says this  was a \textbf{very strong} Phase 1 proposal.
    \item All four criteria received the mark \textbf{6 (Excellent)} on the 1--7 scale. Is this scale correct?
    \item On the Research Council scale, mark 6 means the proposal
          \emph{``successfully addresses essentially all aspects of the criterion and has only minor shortcomings.''}
    \item The path to a competitive Phase 2 application is therefore essentially about
          \textbf{closing the gap from 6 to 7 (Outstanding)} by addressing the specific
          \emph{minor shortcomings} hinted at in each section.
  \end{itemize}
  \vfill
  \begin{block}{Structure of these notes}
    \begin{enumerate}
      \item A less in-depth interpretation of the feedback.
      \item A discussion of explicit things we can address in the full proposal.
    \end{enumerate}
  \end{block}
\end{frame}

%=================================================
\section{Part 1: A less in-depth interpretation}
%=================================================

\begin{frame}{Excellence -- Quality of R\&D activities}
  This section contains the most concrete criticism. Two distinct points are flagged.
  \vfill
  \begin{block}{1. Research questions need sharper articulation}
    Each of the four tracks should open with a \emph{small set of falsifiable scientific questions},
    rather than thematic descriptions.
  \end{block}
  \begin{block}{2. Machine-learning methods need more detail}
    The ML-guided control component is the methodological area the panel found least convincing.
    The full application should include:
    \begin{itemize}
      \item concrete algorithms and model architectures,
      \item training and data strategies,
      \item validation protocols,
      \item the interface between ML and the physics/chemistry tracks.
    \end{itemize}
    A dedicated methods section for ML-guided quantum control is likely warranted.
  \end{block}
\end{frame}

%-------------------------------------------------
\begin{frame}{Impact}
  Two implicit concerns stand out.
  \vfill
  \begin{block}{Scale}
    The panel notes that the center is
    \reviewerquote{relatively small compared with major international centers}
    \vspace{0.3em}
    This is the only comparative remark in the whole report -- and it is not a compliment.
    Phase 2 should either
    \begin{itemize}
      \item grow the scope and team to match international peers, or
      \item explicitly argue why the chosen size delivers \emph{more impact per krone} than larger competitors.
    \end{itemize}
  \end{block}
  \begin{block}{Generic impact narrative}
    The current text stays at the level of \emph{application areas} without committing to specific deliverables,
    pathways, or partners. Phase 2 requires explicit dissemination, exploitation, and open-science
    plans, plus communication and engagement with different target audiences.
  \end{block}
\end{frame}

%-------------------------------------------------
\begin{frame}{Implementation}
  \reviewerquote{some uncertainty exists regarding continuity and transition of leadership}
  \vfill
  Succession uncertainty is one of the recurring reasons SFF panels downgrade otherwise excellent centers.
  \vfill
  \begin{block}{Phase 2 should include a written succession plan}
    \begin{itemize}
      \item Who can step in as director if needed?
      \item How are deputy and co-director roles structured?
      \item How is leadership of each of the four tracks backed up?
      \item How does the center survive the departure of a key PI?
    \end{itemize}
  \end{block}
  \begin{block}{Additional Phase 2 requirements (not yet present)}
    \begin{itemize}
      \item An explicit training environment beyond what individual groups offer.
      \item Appropriate management and governance structures.
      \item A description of the physical organisation supporting cross-track collaboration.
      \item If gender-imbalanced: concrete plans to qualify under-represented-gender researchers
            for senior positions.
    \end{itemize}
  \end{block}
\end{frame}

%=================================================
\section{Part 2: Explicit things we can address}
%=================================================

\begin{frame}{How to read the comments}
  Our goal is to prepare a much larger and more competitive Phase 2 SFF proposal.
  \vfill
  \begin{itemize}
    \item We should read the comments not as criticism, but as indications of where reviewers want
          \textbf{additional depth, precision, and evidence}.
    \item I think that for  highly competitive SFF evaluations, proposals with uniformly excellent scores are
          distinguished by how convincingly they address \emph{these} details.
  \end{itemize}
\end{frame}

%-------------------------------------------------
\begin{frame}{1. The clearest weakness: explicit research questions}
  \reviewerquote{The proposal would benefit from clearer articulation of the specific research questions within each track.}
  \vfill
  The reviewers liked the four-track structure but wanted more than broad themes.
  \vfill
  \begin{columns}[T]
    \column{0.45\textwidth}
    \textbf{Instead of:}
    \begin{itemize}
      \item Quantum control of many-body systems
      \item Machine learning for quantum dynamics
    \end{itemize}

    \column{0.55\textwidth}
    \textbf{Formulate questions such as:}
    \begin{itemize}\small
      \item What are the fundamental limits on controllability in interacting many-body systems?
      \item Can agent-based AI (with reinforcement learning) discover control protocols beyond Pontryagin-optimal solutions?
      \item Can quantum Fisher information provide universal bounds on control fidelity?
      \item Can generative models compress experimentally relevant quantum states?
    \end{itemize}
  \end{columns}
  \vfill
  \begin{center}
    \emph{The best SFF proposals resemble a collection of major open problems, not activities.}
  \end{center}
\end{frame}

%-------------------------------------------------
\begin{frame}{2. Machine learning needs greater methodological precision}
  \reviewerquote{more detailed descriptions of certain methods -- particularly those related to machine-learning-guided control.}
  \vfill
  The panel accepted that ML is important, but were not yet convinced it is sufficiently specified.
  \vfill
  \begin{columns}[T]
    \column{0.45\textwidth}
    \textbf{Name explicit ML methodologies:}
    \begin{itemize}\small
      \item Agent-based AI and Reinforcement learning
      \item Physics-informed neural networks
      \item Bayesian optimization
      \item Diffusion models
      \item Transformer architectures
      \item Generative quantum state models
      \item Graph neural networks
    \end{itemize}

    \column{0.5\textwidth}
    \textbf{For each method, explain:}
    \begin{itemize}\small
      \item why it is needed,
      \item what scientific question it addresses,
      \item why existing methods are insufficient.
    \end{itemize}
  \end{columns}
\end{frame}

%-------------------------------------------------
\begin{frame}{2. ML as a research theme, not just a tool}
  A common weakness in AI-related proposals is that ML appears as a \emph{tool} rather than a
  \emph{research theme}.
  \vfill
  \begin{block}{The proposal should argue}
    New machine-learning methodologies will emerge \textbf{from quantum control challenges themselves}.
  \end{block}
  \vfill
  \begin{block}{Examples of methodological novelty}
    \begin{itemize}
      \item Reinforcement learning for quantum control under uncertainty.
      \item Scientific foundation models for quantum experiments.
      \item Physics-informed generative models.
      \item Hybrid quantum-classical learning algorithms.
    \end{itemize}
  \end{block}
\end{frame}

%-------------------------------------------------
\begin{frame}{3. Clarify scientific ambition versus feasibility}
  The reviewers repeatedly praise the ambition, but Phase 2 should make the relationship
  between risk levels much more explicit.
  \vfill
  \begin{columns}[T]
    \column{0.5\textwidth}
    \textbf{Flagship challenges}
    \begin{itemize}
      \item Fault-tolerant quantum control.
      \item AI-designed quantum sensors.
      \item Autonomous quantum experiments.
    \end{itemize}

    \column{0.5\textwidth}
    \textbf{Fallback pathways}\\
    If a flagship goal fails:
    \begin{itemize}
      \item What is still learned?
      \item What publishable results remain?
    \end{itemize}
  \end{columns}
  \vfill
  \begin{center}
    SFF panels routinely ask: \emph{What happens if the most ambitious ideas fail?}
  \end{center}
\end{frame}

%-------------------------------------------------
\begin{frame}{4. Leadership succession and long-term sustainability}
  \reviewerquote{some uncertainty exists regarding continuity and transition of leadership}
  \vfill
  The panel implicitly asked:
  \begin{itemize}
    \item How dependent is the center on a single person?
    \item What happens after 5--10 years?
    \item How are younger leaders being developed?
  \end{itemize}
  \vfill
  \begin{block}{Explicit leadership pipeline}
    \begin{itemize}
      \item Deputy director(s).
      \item Future work-package leaders.
      \item Early-career leadership positions.
      \item Documented succession plans.
    \end{itemize}
  \end{block}
\end{frame}

%-------------------------------------------------
\begin{frame}{5. Demonstrate stronger center-level synergy}
  The next proposal must repeatedly answer:
  \begin{center}
    \large \emph{Why does this need to be a Center of Excellence, rather than a collection of excellent groups?}
  \end{center}
  \vfill
  \begin{block}{Cross-track flagship projects}
    \begin{itemize}
      \item Quantum sensing $+$ AI.
      \item Quantum materials $+$ machine learning.
      \item Quantum chemistry $+$ control theory.
      \item Experimental platforms $+$ digital twins.
    \end{itemize}
  \end{block}
  \vfill
  \begin{center}
    Each major breakthrough should \textbf{require} collaboration between several tracks.
  \end{center}
\end{frame}

%-------------------------------------------------
\begin{frame}{6. Strengthen international positioning}
  \reviewerquote{relatively small compared with major international centers}
  \vfill
  This reveals a comparison class. For Phase 2, compare explicitly to:
  \vfill
  \begin{columns}[T]
    \column{0.5\textwidth}
    \begin{itemize}
      \item QuICS (Maryland)
      \item IQIM (Caltech)
      \item IQC Waterloo
    \end{itemize}

    \column{0.5\textwidth}
    \begin{itemize}
      \item CQT Singapore
      \item Munich Quantum Valley
      \item Quantum Delta NL
    \end{itemize}
  \end{columns}
  \vfill
  \begin{block}{Then explain}
    \begin{itemize}
      \item Where Norway can \emph{lead}.
      \item What niche is uniquely Norwegian.
      \item Why the center will become internationally indispensable.
    \end{itemize}
  \end{block}
\end{frame}

%-------------------------------------------------
\begin{frame}{7. Expand the impact narrative}
  The current impact section is strong but generic (secure communication, advanced materials,
  healthcare, clean energy). Phase 2 should make this much more concrete.
  \vfill
  \begin{columns}[T]
    \column{0.33\textwidth}
    \textbf{Scientific impact}
    \begin{itemize}\small
      \item New theory of quantum control.
      \item New algorithms.
      \item Open-source software.
    \end{itemize}

    \column{0.33\textwidth}
    \textbf{Technology impact}
    \begin{itemize}\small
      \item Quantum sensors.
      \item Quantum materials discovery.
      \item Quantum computing architectures.
    \end{itemize}

    \column{0.33\textwidth}
    \textbf{Human capital impact}
    \begin{itemize}\small
      \item Number of PhDs.
      \item Industrial internships.
      \item International mobility.
    \end{itemize}
  \end{columns}
\end{frame}

%-------------------------------------------------
\begin{frame}{8. What is missing entirely from the criticism?}
  Interestingly, there are \textbf{no} criticisms regarding:
  \begin{itemize}
    \item scientific quality,
    \item leadership competence,
    \item publication record,
    \item feasibility,
    \item interdisciplinarity.
  \end{itemize}
  \vfill
  \begin{center}
    \emph{This is very encouraging.} The panel is convinced the team can do the work.
  \end{center}
  \vfill
  \begin{block}{The next proposal should focus less on proving excellence and more on:}
    \begin{enumerate}
      \item Sharper scientific questions.
      \item More detailed methodology.
      \item Clearer ML strategy.
      \item Stronger center-level integration.
      \item Leadership succession.
      \item A stronger international vision.
    \end{enumerate}
  \end{block}
\end{frame}

%=================================================
\section{Overall interpretation}
%=================================================

\begin{frame}{Reading between the lines}
  \reviewerquote{We are convinced that this is an excellent team pursuing an important and
  potentially transformative topic. For the full SFF proposal we would like to see a more
  detailed scientific roadmap, especially regarding machine-learning-guided quantum control,
  clearer articulation of the key scientific questions, and stronger evidence that the center
  will function as more than the sum of its individual research groups.}
  \vfill
  \begin{itemize}
    \item A very favourable position going into a second-round SFF application.
    \item The reviewers did not identify fundamental weaknesses.
    \item They asked for \textbf{greater precision}, \textbf{sharper focus}, and a
          \textbf{more developed center strategy}.
  \end{itemize}
\end{frame}

%-------------------------------------------------
\begin{frame}{Summary: priorities for the Phase 2 proposal}
  \begin{enumerate}
    \item \textbf{Leadership continuity and succession plan} -- the only explicit Implementation concern.
    \item \textbf{Detailed methodology for ML-guided quantum control} -- dedicated methods section.
    \item \textbf{Sharpened research questions} per track -- falsifiable, ambitious, specific.
    \item \textbf{Scale and international positioning} -- either grow, or justify the chosen scale against named peer centers.
    \item \textbf{Center-level synergy} -- cross-track flagship projects that no single group could deliver.
    \item \textbf{Risk management} -- flagship goals with explicit fallback pathways.
    \item \textbf{Phase 2 requirements} -- dissemination, exploitation, open science, training environment, gender plans, governance.
  \end{enumerate}
\end{frame}


\end{document}
